% GENERATED FILE. Edit canonical lessons, then run npm run generate:books.
% canonical-source-hash: fnv1a64:6025b54cf375a307
% canonical-lessons: RU-C295-pozhar, RU-C295-dym, RU-C295-goret, RU-C295-zagoratsya, RU-C295-ognetushitel

\chapter{Fire}
\label{ch:ru-fire}
\hlchaptermodality{295}

\begin{chapteropening}
\textbf{By the end of this chapter:} \emph{I can talk about a fire: a blaze, smoke, burning, catching fire and a fire extinguisher.}

Complete the last lesson of chapter 295: I can talk about a fire: a blaze, smoke, burning, catching fire and a fire extinguisher.
\end{chapteropening}

\section[\texorpdfstring{\ru{пожар} (pozhár)}{пожар (pozhár)}]{\ru{пожар} (pozhár) — a fire (a blaze)}

\label{lesson:RU-C295-pozhar}



\begin{quote}
Before the new one: say the Russian for to confess, then the Russian for to acquit.
\end{quote}

\subsection*{You'll want to know: \ru{пожар}}
\textbf{\ru{пожар}} (\emph{pozhár}) — \textquotedblleft{}a fire (a blaze)\textquotedblright{}.

\begin{quote}
\ru{В} \ru{доме} \ru{напротив} \ru{был} \ru{пожар}.
\end{quote}

\begin{quote}
\emph{V dóme naprótiv byl pozhár.}
\end{quote}

\begin{quote}
There was a fire in the building opposite.
\end{quote}

A fire that destroys. \emph{\ru{Пожарный}} — a firefighter.

\begin{cousinweb}
\emph{\ru{пожар}} is built on \emph{\ru{по}-} and \emph{\ru{жар}} \textquotedblleft{}heat\textquotedblright{}, from the same Slavic root as \emph{\ru{гореть}} \textquotedblleft{}to burn\textquotedblright{}; EtymDB lists \emph{\ru{пожар}} with no parent.
\end{cousinweb}

\begin{grammarlens}[title={What you've built}]
One more word for telling, explaining and describing at B1.
\end{grammarlens}

\subsection*{Your turn}
\begin{itemize}
\raggedright
  \item \emph{Say it:} \emph{pozhár}
  \item \emph{Say it:} \emph{pozhár}, once more
  \item \emph{Say it:} a sentence of your own with \emph{pozhár}
  \item \emph{Recall:} say the Russian for to confess, then the Russian for to acquit, then say \emph{pozhár} again
\end{itemize}

\begin{tcolorbox}[breakable,colback=teal!4,colframe=teal!35!black,title={Before you move on}]
What does \ru{пожар} mean? (\textquotedblleft{}A fire (a blaze)\textquotedblright{}.) Say it once more.
\end{tcolorbox}
\section[\texorpdfstring{\ru{дым} (dym)}{дым (dym)}]{\ru{дым} (dym) — smoke}

\label{lesson:RU-C295-dym}



\begin{quote}
Before the new one: say the Russian for to acquit, then the Russian for a fire.
\end{quote}

\subsection*{You'll want to know: \ru{дым}}
\textbf{\ru{дым}} (\emph{dym}) — \textquotedblleft{}smoke\textquotedblright{}.

\begin{quote}
\ru{Из} \ru{окна} \ru{шёл} \ru{чёрный} \ru{дым}.
\end{quote}

\begin{quote}
\emph{Iz okná shyol chyórnyy dym.}
\end{quote}

\begin{quote}
Black smoke was coming out of the window.
\end{quote}

\emph{\ru{Нет} \ru{дыма} \ru{без} \ru{огня}} — there is no smoke without fire.

\begin{cousinweb}
\emph{\ru{дым}} is inherited from Proto-Slavic \emph{dymŭ} \textquotedblleft{}smoke\textquotedblright{}, from PIE \emph{d\textsuperscript{h}uh\textsubscript{2}mós} \textquotedblleft{}smoke\textquotedblright{} (EtymDB); Latin \emph{fūmus} \textquotedblleft{}smoke\textquotedblright{} is a cousin.
\end{cousinweb}

\begin{grammarlens}[title={What you've built}]
One more word for telling, explaining and describing at B1.
\end{grammarlens}

\subsection*{Your turn}
\begin{itemize}
\raggedright
  \item \emph{Say it:} \emph{dym}
  \item \emph{Say it:} \emph{dym}, once more
  \item \emph{Say it:} a sentence of your own with \emph{dym}
  \item \emph{Recall:} say the Russian for to acquit, then the Russian for a fire, then say \emph{dym} again
\end{itemize}

\begin{tcolorbox}[breakable,colback=teal!4,colframe=teal!35!black,title={Before you move on}]
What does \ru{дым} mean? (\textquotedblleft{}Smoke\textquotedblright{}.) Say it once more.
\end{tcolorbox}
\section[\texorpdfstring{\ru{гореть} (gorét)}{гореть (gorét)}]{\ru{гореть} (gorét) — to burn, to be on fire}

\label{lesson:RU-C295-goret}



\begin{quote}
Before the new one: say the Russian for a fire, then the Russian for smoke.
\end{quote}

\subsection*{You'll want to know: \ru{гореть}}
\textbf{\ru{гореть}} (\emph{gorét}) — \textquotedblleft{}to burn, to be on fire\textquotedblright{}.

\begin{quote}
\ru{В} \ru{лесу} \ru{горела} \ru{трава}.
\end{quote}

\begin{quote}
\emph{V lesú goréla travá.}
\end{quote}

\begin{quote}
Grass was burning in the forest.
\end{quote}

Imperfective; present \emph{\ru{горю}}, \emph{\ru{горишь}}. Perfective \emph{\ru{сгореть}} (\emph{sgorét}) — to burn down. Of a light: \emph{\ru{свет} \ru{горит}} — the light is on.

\begin{cousinweb}
\emph{\ru{гореть}} is inherited from Proto-Slavic \emph{gorěti} \textquotedblleft{}to burn\textquotedblright{}, which EtymDB traces to PIE \emph{g\textsuperscript{w}\textsuperscript{h}er-} \textquotedblleft{}warm, hot\textquotedblright{}.
\end{cousinweb}

\begin{grammarlens}[title={What you've built}]
One more word for telling, explaining and describing at B1.
\end{grammarlens}

\subsection*{Your turn}
\begin{itemize}
\raggedright
  \item \emph{Say it:} \emph{gorét}
  \item \emph{Say it:} \emph{gorét}, once more
  \item \emph{Say it:} a sentence of your own with \emph{gorét}
  \item \emph{Recall:} say the Russian for a fire, then the Russian for smoke, then say \emph{gorét} again
\end{itemize}

\begin{tcolorbox}[breakable,colback=teal!4,colframe=teal!35!black,title={Before you move on}]
What does \ru{гореть} mean? (\textquotedblleft{}To burn, to be on fire\textquotedblright{}.) Say it once more.
\end{tcolorbox}
\section[\texorpdfstring{\ru{загораться} (zagorátsya)}{загораться (zagorátsya)}]{\ru{загораться} (zagorátsya) — to catch fire}

\label{lesson:RU-C295-zagoratsya}



\begin{quote}
Before the new one: say the Russian for smoke, then the Russian for to burn.
\end{quote}

\subsection*{You'll want to know: \ru{загораться}}
\textbf{\ru{загораться}} (\emph{zagorátsya}) — \textquotedblleft{}to catch fire\textquotedblright{}.

\begin{quote}
\ru{От} \ru{свечи} \ru{загорелась} \ru{штора}.
\end{quote}

\begin{quote}
\emph{Ot svechí zagorélas shtóra.}
\end{quote}

\begin{quote}
The curtain caught fire from a candle.
\end{quote}

Imperfective; perfective \emph{\ru{загореться}} (\emph{zagorétsya}). Also of lights coming on.

\begin{cousinweb}
\emph{\ru{загораться}} is built on \emph{\ru{за}-}, which marks a start, and \emph{\ru{гореть}} \textquotedblleft{}to burn\textquotedblright{}, inherited from Proto-Slavic \emph{gorěti} (EtymDB): to start burning.
\end{cousinweb}

\begin{grammarlens}[title={What you've built}]
One more word for telling, explaining and describing at B1.
\end{grammarlens}

\subsection*{Your turn}
\begin{itemize}
\raggedright
  \item \emph{Say it:} \emph{zagorátsya}
  \item \emph{Say it:} \emph{zagorátsya}, once more
  \item \emph{Say it:} a sentence of your own with \emph{zagorátsya}
  \item \emph{Recall:} say the Russian for smoke, then the Russian for to burn, then say \emph{zagorátsya} again
\end{itemize}

\begin{tcolorbox}[breakable,colback=teal!4,colframe=teal!35!black,title={Before you move on}]
What does \ru{загораться} mean? (\textquotedblleft{}To catch fire\textquotedblright{}.) Say it once more.
\end{tcolorbox}
\section[\texorpdfstring{\ru{огнетушитель} (ognetushítel)}{огнетушитель (ognetushítel)}]{\ru{огнетушитель} (ognetushítel) — a fire extinguisher}

\label{lesson:RU-C295-ognetushitel}



\begin{quote}
Before the new one: say the Russian for to burn, then the Russian for to catch fire.
\end{quote}

\subsection*{You'll want to know: \ru{огнетушитель}}
\textbf{\ru{огнетушитель}} (\emph{ognetushítel}) — \textquotedblleft{}a fire extinguisher\textquotedblright{}.

\begin{quote}
\ru{Огнетушитель} \ru{висит} \ru{на} \ru{стене}.
\end{quote}

\begin{quote}
\emph{Ognetushítel visít na stené.}
\end{quote}

\begin{quote}
The fire extinguisher hangs on the wall.
\end{quote}

Masculine; genitive \emph{\ru{огнетушителя}}.

\begin{cousinweb}
\emph{\ru{огнетушитель}} joins \emph{\ru{огонь}} \textquotedblleft{}fire\textquotedblright{}, inherited through Old East Slavic \emph{ognĭ} \textquotedblleft{}fire\textquotedblright{} (EtymDB), and \emph{\ru{тушить}} \textquotedblleft{}to put out\textquotedblright{}: a fire-quencher.
\end{cousinweb}

\begin{grammarlens}[title={What you've built}]
One more word for telling, explaining and describing at B1.
\end{grammarlens}

\subsection*{Your turn}
\begin{itemize}
\raggedright
  \item \emph{Say it:} \emph{ognetushítel}
  \item \emph{Say it:} \emph{ognetushítel}, once more
  \item \emph{Say it:} a sentence of your own with \emph{ognetushítel}
  \item \emph{Recall:} say the Russian for to burn, then the Russian for to catch fire, then say \emph{ognetushítel} again
\end{itemize}

\begin{tcolorbox}[breakable,colback=teal!4,colframe=teal!35!black,title={Before you move on}]
What does \ru{огнетушитель} mean? (\textquotedblleft{}A fire extinguisher\textquotedblright{}.) Say it once more.
\end{tcolorbox}
